\paragraph{Endpoint pushforward lemma for the relative-wave argument.}
Use standard Lorentz coordinates, with source
$p(u)=(\cosh u,\sinh u,0)$ and an unoriented target geodesic
represented by a unit spacelike normal
$n=(N_0,N_1,N_2)$, so that
\[
 N_0^2-N_1^2-N_2^2=-1.
\]
Put $a=17$, $B=12\sqrt2$, $R=\operatorname{arcosh}a$ and
$\tau(u)=\operatorname{arsinh}(B\cosh u)$.
Fix a compact interval $J=[\ell_-,\ell_+]\subset(0,R)$.
On a choice of sign of $n$, let
\[
 r_n(u)=\operatorname{arsinh}(N_0\cosh u-N_1\sinh u),
 \qquad \ell=\tau(u)-r_n(u).
\]
Only the part with $r_n>0$ is considered for this phase.  On its
preimage of $J$, this positivity is automatic because
$r_n=\tau-\ell\geq R-\ell_+>0$.
For a fixed smooth function $\chi$ supported in the interior of $J$,
push forward the principal measure
\[
 \chi(\ell)\frac{a}{2\cosh^{3/2}\tau(u)\cosh^{1/2}r_n(u)}\,du
\]
by the map $u\mapsto\ell$.  There are constants $M_J,C_J>0$,
independent of the normal, such that if $|N_1|\geq M_J$, this
pushforward is $h_n(\ell)\,d\ell$ and its extension by zero satisfies
\begin{equation}\label{A:endpoint-BV-single}
 \|h_n\|_{L^1(\mathbb R)}+
    \operatorname{Var}_{\mathbb R}(h_n)
 \leq \frac{C_J}{|N_1|}.
\end{equation}
The same conclusion holds after multiplying the displayed measure
by any factor whose expression in $(v,\ell)$ below has a uniformly
bounded Lipschitz norm on $[-1,1]\times J$.

\begin{proof}
Set $v=\tanh u$ and define
\[
 D(v)=\sqrt{a^2-v^2},\qquad
 E(v,\ell)=D(v)\cosh\ell-B\sinh\ell.
\]
The identities
\begin{align*}
 \cosh\tau&=\frac{D(v)}{\sqrt{1-v^2}},&
 \cosh r_n&=\frac{E(v,\ell)}{\sqrt{1-v^2}},\\
 F(v,\ell)&=N_0-N_1v-B\cosh\ell+D(v)\sinh\ell=0
\end{align*}
follow from $r_n=\tau-\ell$.  Conversely, on $J$ the latter
equation implies this equality with $r_n>0$, since
\[
 B\cosh\ell-D\sinh\ell
 \geq B\cosh\ell_+-a\sinh\ell_+>0.
\]
The elementary uniform bounds are
\[
 B\leq D\leq a,\qquad
 B e^{-\ell_+}\leq E\leq a\cosh\ell_+ +B\sinh\ell_+.
\]
The measure simplifies exactly, including the Jacobian:
\[
 \frac{a}{2\cosh^{3/2}\tau\cosh^{1/2}r_n}\,du
 =\frac{a}{2D^{3/2}E^{1/2}}\,dv.
\]
Moreover
\[
 F_v=-N_1-\frac{v\sinh\ell}{D},\qquad
 F_\ell=E,\qquad
 \frac{dv}{d\ell}=
 \frac{E}{N_1+v\sinh\ell/D}.
\]
Choose
$M_J=\max\{1,2\sinh(\ell_+)/B\}$.
For $|N_1|\geq M_J$, $F_v$ has constant nonzero sign on
$[-1,1]\times J$.  There is at most one root $v=v_n(\ell)$
for each $\ell\in J$, and on the interval where it exists
the density is exactly
\begin{equation}\label{A:endpoint-density-exact}
 h_n(\ell)=\chi(\ell)
 \frac{a\sqrt{E(v_n(\ell),\ell)}}
 {2D(v_n(\ell))^{3/2}
  |N_1+v_n(\ell)\sinh\ell/D(v_n(\ell))|}.
\end{equation}
The allowed set of $\ell$ is an interval, possibly empty:
at $v=\pm1$,
\[
 F(\pm1,\ell)=N_0\mp N_1-B e^{-\ell},
\]
and each endpoint condition is monotone in $\ell$.
These are precisely the possible limiting Busemann endpoints;
no separation assumption on different normals' endpoints is needed.

Here are explicit bounds sufficient to differentiate the density.
Write $Z=N_1+v\sinh\ell/D$ and
$E_1=a\cosh\ell_+ +B\sinh\ell_+$.
Then $|Z|\geq|N_1|/2$ and
$|v_n'|\leq2E_1/|N_1|$.
Since $D_v=-v/D$ and $(v/D)_v=a^2/D^3$, differentiation
along the root gives
\[
 |D'|\leq\frac{2E_1}{B|N_1|},\qquad
 |Z'|\leq\frac{2a^2E_1\sinh\ell_+}{B^3|N_1|}
                  +\frac{\cosh\ell_+}{B}.
\]
Also $|E'|$ is bounded by a constant depending only on $J,a,B$,
and $E$ is bounded away from zero.  Applying these inequalities
to \eqref{A:endpoint-density-exact} proves
$\sup|h_n|+\sup|h_n'|\leq C_J/|N_1|$ in its interval of
definition.  At most two endpoints occur; extending by zero adds
jumps of total variation at most $2\sup|h_n|$.
This proves \eqref{A:endpoint-BV-single}.  The product rule proves
the assertion with a uniformly bounded Lipschitz factor.
\end{proof}

\paragraph{Summability over large-height axes and all their translates.}
For $|N_2|>1$ choose the sign of the normal with $N_0>0$ and put
\[
 c=N_0-|N_1|>0,\qquad A_2=\sqrt{N_2^2-1}.
\]
If the phase has a value in $J$, its defining equation gives
\[
 c\leq N_0-N_1v
       =B\cosh\ell-D\sinh\ell
       \leq B e^{-\ell}\leq C_0,
 \qquad C_0=B e^{-\ell_-}.
\]
The normal equation implies
\[
 |N_1|=\frac{N_2^2-1-c^2}{2c}.
\]
Thus, whenever $N_2^2\geq2(1+C_0^2)$,
\[
 |N_1|\geq\frac{N_2^2}{4c},\qquad
 \|h_n\|_{BV}\leq\frac{4C_Jc}{N_2^2},
\]
after increasing the fixed height cutoff if necessary to ensure
$|N_1|\geq M_J$.

The source stabilizer is a boost of length $L_0>0$.
Write $N_0=A_2\cosh u_0$ and $N_1=A_2\sinh u_0$.
Its translates have
\[
 c_m=A_2e^{-|u_0+mL_0|},\qquad m\in\mathbb Z.
\]
The terms with $c_m\leq C_0$ consist of two tails of a geometric
progression (or, if the cutoff exceeds their maximum, the whole
progression).  Splitting at a nearest integer to $-u_0/L_0$ gives
the uniform estimate
\[
 \sum_{m:\,c_m\leq C_0}c_m
 \leq \frac{2C_0}{1-e^{-L_0}}.
\]
Consequently each reduced target axis of height $|N_2|$ contributes
at most $C/N_2^2$ to the total BV norm of all its source translates.
Statement A.3.1, with its double-coset bookkeeping, gives
$O(Y)$ reduced axes of height at most $Y$; in our coordinates
$|N_2|=3|z|$ for the disjoint target axes.  Dyadic summation now gives
\[
 \sum_{\substack{\text{reduced axes}\ |N_2|\geq Y_0}}
 \ \sum_{m\in\mathbb Z}\|h_{n,m}\|_{BV}
 \ll\sum_{j\geq0}\frac{2^jY_0}{(2^jY_0)^2}<\infty.
\]
No assertion about spacing of Busemann endpoints is used.

\paragraph{Finite height and intersecting axes.}
There are finitely many reduced axes below $Y_0$, including the
finite set of intersecting-axis double cosets.  For each such axis,
the boost formulas express $N_1(m)$ as a linear combination of
$e^{mL_0}$ and $e^{-mL_0}$ with both coefficients nonzero, unless
the axis shares an endpoint with the source.  The latter case has
already been excluded in A.3.1 by discreteness and cocompactness;
equality of the axes is separately excluded by the integral normals.
It follows that $|N_1(m)|\asymp e^{|m|L_0}$ as $m\to\pm\infty$,
with constants allowed to depend on this finite list of axes.
Therefore
\[
 \sum_{m:\,|N_1(m)|\geq M_J}|N_1(m)|^{-1}<\infty.
\]
The preceding single-axis BV estimate applies to these tails,
with either sign of the normal as needed.  In particular this
treats $N_2=0$ (the intersecting Pell families), without using
the large-height formula involving $c$.
Only finitely many axes/translates with $|N_1|<M_J$ remain.
Their interior stationary points must be treated separately by
ordinary stationary phase; they are not absorbed into the endpoint
estimate.  A normal with $|N_2|=1$ would be asymptotic to the source
and is excluded; in the present integral coordinates it cannot
occur in any event, since $N_2=3z$.

\paragraph{Uniform outgoing amplitudes.}
The same summability holds for the exact high-frequency amplitude
in the outgoing relative-wave factors.  To specify the needed
regularity, the outgoing Volterra estimates in A.4.8 imply, after
removing $e^{-itr}$, a bounded amplitude $b_t(r)$ with
\[
 |b_t(r)|\leq C,\qquad |\partial_r b_t(r)|\leq Ce^{-2r},
 \qquad t\geq1,\quad r\geq r_0>0.
\]
For the normal derivative factor this follows as well: if
$w_t(r)=\sqrt2e^{itr}\sqrt{\cosh r}\,f_{1/2+it}(r)$,
then $w_t=1+O(t^{-1}e^{-2r})$, $w_t'=O(e^{-2r})$, and
$w_t''=2itw_t'+Vw_t$; substituting in $f_s'$ gives the same
uniform bound for the derivative amplitude.  Scalar gamma factors,
after their leading powers and phases have been removed, are
bounded independently of $r$ and cause no difficulty.
On the phase surface,
\[
 \cosh\tau=\frac{D}{\sqrt{1-v^2}},\qquad
 r_n=\tau-\ell,\qquad
 \cosh r_n=\frac{E}{\sqrt{1-v^2}}.
\]
Thus $e^{-2\tau}+e^{-2r_n}\ll_J1-v^2$, whereas
$|\partial_v\tau|+|\partial_v r_n|\ll_J(1-v^2)^{-1}$
and $\partial_\ell r_n=-1$ at fixed $v$.
The product $b_t(\tau)b_t(r_n)$ consequently has bounded first
derivatives on $(-1,1)\times J$ and a uniformly bounded Lipschitz
extension to $[-1,1]\times J$, uniformly in $t\geq1$.
The factor version of the lemma therefore applies to the complete
outgoing symbol, not only its principal term.

\paragraph{Oscillatory consequence and precise scope.}
The sum of all these large-$|N_1|$ pushforwards is a compactly
supported BV density $h_t$ with norm bounded uniformly for $t\geq1$.
Integration by parts in the distributional sense yields
\[
 \left|\int_{\mathbb R}e^{it\ell}h_t(\ell)\,d\ell\right|
 \leq\frac{\operatorname{Var}(h_t)}{t}\ll_J t^{-1}.
\]
This is smaller than the $t^{-1/2}$ contribution of a nondegenerate
interior stationary point in the same oscillatory integral.
The lemma proves the endpoint comparison needed in a separately
normalized relative-wave argument.  It does not itself identify
the finite stationary set, prove its leading coefficient nonzero,
or assert an $\Omega$ theorem.
